%% ----------------------------------------------------------------
%% Probability problem sheet (lectures: probability, experimentalDesign, bayes)
%% ----------------------------------------------------------------
\documentclass{sotonExamBoxes}    % use option [answers] to include them
%% ----------------------------------------------------------------
\usepackage{graphicx}
\usepackage{alltt}
\usepackage{bm}
\usepackage{amssymb}
\usepackage{amsfonts}
\usepackage{amsmath}
\thinmuskip=1mu
%% macros.tex -- shared macro definitions for the problem sheets.
%%
%% Every sheet inputs this instead of carrying its own copy. Input it after
%% the \usepackage lines: the definitions below use bm, amsmath and amssymb.
%%
%% Six macros were defined inconsistently across the sheets when this file
%% was made; the choices are marked CONFLICT below.




\usepackage{stmaryrd}

\renewcommand{\ttdefault}{pcr}

\newcommand{\tr}{\textsf{T}}

\newcommand{\e}[1]{{\rm e}^{#1}}

\newcommand{\E}[1]{{\rm e}^{{\displaystyle #1}}}

\newcommand{\logg}[1]{\log\left(#1\right)}

\newcommand{\Prob}[1]{\mathbb{P}\left(#1\right)}

\newcommand{\Step}{\mathop{\mathrm{O}\hspace{-15pt}\rule[8.5pt]{11pt}{1.2pt}\hspace{2pt}}}

\newcommand\diag{\mathop{\mathrm{diag}}}

\newcommand{\grad}{\bm{\nabla}}

\newcommand{\dd}{\mathrm{d}}

\newcommand{\class}{\mathcal{C}}

\newcommand{\data}{\mathcal{D}}

%% CONFLICT: gaussianProcess.tex had {U}{eur}{b}{n} (Euler bold);
%% every other sheet had this. Matrices in that one sheet now look
%% different -- see the report.

\DeclareMathAlphabet{\mat}{OT1}{cmss}{bx}{n}

%% The identity needs to be distinguishable. In cmss a bold I is a plain
%% vertical bar with no crosspieces, so \mat{I} reads as a rule rather than a
%% letter. Euler's bold I keeps its serifs, so the identity alone is set in
%% Euler while everything else stays cmss.
\DeclareMathAlphabet{\matEuler}{U}{eur}{b}{n}
\newcommand{\Id}{\matEuler{I}}

%% CONFLICT and BUG FIX: 16 sheets defined this using \qpartial,
%% which is defined nowhere, so \pd was broken in all of them.
%% Only problem2.tex had the working \partial form, kept here.

\newcommand{\pd}[2]{\ensuremath{\frac{\partial #1}{\partial #2}}\xspace}

\newcommand{\M}[1]{\ensuremath{\boldsymbol{#1}}\xspace}

\newcommand{\bst}{\rule{0mm}{5mm}}

%% CONFLICT: some sheets omitted the \! . Kept the tighter spacing,
%% which is what every sheet that actually uses \av had.

\newcommand{\av}[2][\,]{\mathbb{E}_{#1}\!\left[ {\strut #2} \right]}

\newcommand{\len}[1]{\| #1 \|}

\newcommand{\inner}[2]{\left\langle #1, #2\right\rangle}

\newcommand{\squeeze}{\setlength{\itemsep}{-2pt}}

%% CONFLICT: two spacings, identical output in maths mode.

\newcommand{\bra}[1]{\left( #1 \right)}

\newcommand{\pred}[1]{\left\llbracket { \small #1} \right\rrbracket}

\newcommand{\mln}{\hat{f}_n(\bm{x})}


\newcommand{\Tr}{\mathop{\mathrm{tr}}\,}

%% NOT SHARED: gaussianProcess.tex redefines \det to take an argument
%% (\det{\mat{V}}), while calculus.tex and vectorSpaces.tex use the standard
%% no-argument operator (\det(\mat{X})). The two cannot coexist, so that
%% redefinition stays local to gaussianProcess.tex.

%% NOT SHARED: \ml and \mm mean different things in different sheets.
%% amgm and convexity use a no-argument \ml and \mm = \hat{m}_c(x);
%% ensembleLearning uses a one-argument \ml and \mm = \hat{m}(x).
%% Both forms are in use, so each sheet keeps its own.

\graphicspath{{figures/}}
\usepackage{parskip}
\AddToHook{env/samepage/begin}{%
\AddToHook{para/before}{\nopagebreak}%
\AddToHook{para/after}{\nopagebreak}%
}
\begin{document}
\unitTitle{Probability Problem Sheet}
\authors{Adam Pr\"ugel-Bennett}
\unitcode{AICE3001/AICE6002}
\semester{2}
\year{2022}
\duration{45}{0}

\maketitle

This sheet covers the lectures on probability, on interpreting
experimental results, and on Bayesian inference.  Question~1 practises
the rules of probability, Question~2 error bars and significance tests,
and Question~3 Bayesian updating.  Each ends with a short Python part
(\texttt{import numpy as np} is assumed).  Question~4 is harder and
optional: it does Bayesian updating for normal distributions and for
linear regression.

% ***************************** question 1 **************************

\begin{question}{20}
  \begin{qparts}

    \qpart[4]{Events $A$ and $B$ have $\Prob{A}=0.3$, $\Prob{B}=0.5$
      and $\Prob{A\cap B}=0.1$.  Find $\Prob{A\cup B}$,
      $\Prob{\neg A}$ and $\Prob{A\mid B}$.  Are $A$ and $B$
      independent?
      \begin{answer}
        Adding $\Prob{A}$ and $\Prob{B}$ counts $A\cap B$ twice, so
        $\Prob{A\cup B} = 0.3+0.5-0.1 = 0.7$.  $\Prob{\neg A} = 1-0.3 =
        0.7$.  $\Prob{A\mid B} = \Prob{A\cap B}/\Prob{B} = 0.1/0.5 =
        0.2$.  They are not independent: $\Prob{A}\,\Prob{B} = 0.15 \neq
        0.1$ (equivalently $\Prob{A\mid B} = 0.2 \neq \Prob{A}$).
      \end{answer}}{\lines{5}}

    \qpart[5]{The joint distribution of the weather $W$ and whether
      I carry an umbrella $U$ is
      \begin{center}
        \begin{tabular}{l|cc}
          & $U=$ yes & $U=$ no \\ \hline
          $W=$ sun & 0.05 & 0.55 \\
          $W=$ rain & 0.25 & 0.15
        \end{tabular}
      \end{center}
      Find $\Prob{U=\text{yes}}$, $\Prob{W=\text{rain}\mid
      U=\text{yes}}$ and $\Prob{U=\text{yes}\mid W=\text{rain}}$.
      Comment on the last two.
      \begin{answer}
        Marginalising, $\Prob{U=\text{yes}} = 0.05+0.25 = 0.30$ and
        $\Prob{W=\text{rain}} = 0.25+0.15 = 0.40$.  Then
        $\Prob{\text{rain}\mid\text{yes}} = 0.25/0.30 \approx 0.83$ and
        $\Prob{\text{yes}\mid\text{rain}} = 0.25/0.40 = 0.625$.  The two
        conditional probabilities are different: conditioning is not
        symmetric, and the umbrella does not cause the rain.
      \end{answer}}{\lines{6}}

    \qpart[4]{$X$ is the score on a fair dice.  Find $\av{X}$,
      $\av{X^2}$ and $\mathrm{Var}(X)$.  Use linearity of expectation
      to find $\av{2X+3}$, and find $\mathrm{Var}(2X+3)$.
      \begin{answer}
        $\av{X} = \frac16(1+\cdots+6) = 3.5$ and
        $\av{X^2} = \frac16(1+4+\cdots+36) = \frac{91}{6}$, so
        $\mathrm{Var}(X) = \frac{91}{6}-3.5^2 = \frac{35}{12}$.  Then
        $\av{2X+3} = 2\times3.5+3 = 10$, and since adding a constant does
        not change the spread, $\mathrm{Var}(2X+3) = 4\,\mathrm{Var}(X) =
        \frac{35}{3}$.
      \end{answer}}{\lines{5}}

    \qpart[4]{$X$ is uniform on $(0,1)$ and $Y = -\log(X)/\lambda$ with
      $\lambda>0$.  Use the change of variables formula
      $f_Y(y) = f_X(x)\,|\dd x/\dd y|$ to find the density of $Y$.
      \begin{answer}
        $f_X(x) = 1$ on $(0,1)$, and $x = \e{-\lambda y}$ so
        $|\dd x/\dd y| = \lambda\,\e{-\lambda y}$.  Hence $f_Y(y) =
        \lambda\,\e{-\lambda y}$ for $y>0$: the exponential distribution.
        This is how computers turn uniform random numbers into
        exponential ones.
      \end{answer}}{\lines{5}}

    \qpart[3]{In Python, draw $100\,000$ samples of $Y$ from part (d)
      with $\lambda=2$.  Check that their mean is close to $1/\lambda$,
      and estimate $\Prob{Y\leq 0.5}$ as the average of the indicator
      $[\![Y\leq0.5]\!]$.  Compare with the exact value
      $1-\e{-\lambda/2}$.
      \begin{answer}
\begin{alltt}
rng = np.random.default\_rng(0)\\
lam = 2.0\\
y = -np.log(rng.uniform(size=100000)) / lam\\
print(y.mean(), 1/lam)\\
print(np.mean(y <= 0.5), 1 - np.exp(-lam*0.5))
\end{alltt}
        The mean is $0.500$ and the estimate $0.633$, against $0.632$:
        a probability is the expectation of an indicator.
      \end{answer}}{\lines{5}}

  \end{qparts}
\end{question}

% ***************************** question 2 **************************

\begin{question}{20}
  \begin{qparts}

    \qpart[3]{A classifier gets 1700 of 2000 test examples right.
      Estimate its accuracy and standard error, and write the result
      down sensibly.
      \begin{answer}
        $\hat{p} = 0.85$ and $\Delta = \sqrt{0.85\times0.15/2000}
        \approx 0.008$, so the accuracy is $0.850\pm0.008$ (or
        $(85.0\pm0.8)\%$).
      \end{answer}}{\lines{4}}

    \qpart[4]{A second classifier gets $0.866\pm0.008$ on a
      \emph{different} test set of 2000 examples.  Is it significantly
      better at the 5\% level?
      \begin{answer}
        The gap is $0.016$.  The two estimates are independent, so the
        standard error of the difference is $\sqrt{2}\times0.008 \approx
        0.0113$, not $0.008$.  Then $t\approx1.42$ and, with so many
        examples, $p\approx0.16$: not significant.  The error bars are
        about two standard errors apart, which is not enough.
      \end{answer}}{\lines{5}}

    \qpart[5]{Now both classifiers are run on the \emph{same} 2000
      examples.  The first is right and the second wrong on 20 examples;
      the second is right and the first wrong on 52; they agree on the
      rest.  Let $D_i$ be $+1$, $-1$ or $0$ as the second classifier does
      better, worse or the same on example $i$.  Use a paired t-test on
      the $D_i$.  Why does the answer differ from part (b)?
      \begin{answer}
        $\bar{D} = (52-20)/2000 = 0.016$, the same gap as before.  The
        sample variance of the $D_i$ is about $72/2000 - 0.016^2 \approx
        0.0357$, so the standard error is $\sqrt{0.0357/2000} \approx
        0.0042$ and $t\approx3.8$, giving $p\approx0.0002$: highly
        significant.  On the same examples most of the variation (easy
        and hard examples) is shared and cancels in the differences; only
        the 72 examples where they disagree carry information.
      \end{answer}}{\lines{7}}

    \qpart[3]{You try 20 hyperparameter settings, none of which is
      really any better than the baseline, and test each at the 5\%
      level.  What is the probability that at least one looks
      significant?
      \begin{answer}
        $1-0.95^{20}\approx0.64$.  Running many tests and reporting the
        best is cheating, even if it is done without meaning to.
      \end{answer}}{\lines{3}}

    \qpart[5]{Check part (b) in Python by simulating the null hypothesis:
      two classifiers that are both right with probability $0.85$, each
      tested on its own 2000 examples.  Repeat $100\,000$ times and find
      the fraction of trials in which the accuracies differ by at least
      $0.016$.
      \begin{answer}
\begin{alltt}
rng = np.random.default\_rng(0)\\
a = rng.binomial(2000, 0.85, size=100000) / 2000\\
b = rng.binomial(2000, 0.85, size=100000) / 2000\\
print(np.mean(np.abs(a - b) >= 0.016))
\end{alltt}
        This gives about $0.16$, agreeing with part (b) (the small
        excess over $0.157$ is because the accuracies can only take
        multiples of $1/2000$).  Simulating the null hypothesis works
        even when no formula is available.
      \end{answer}}{\lines{6}}

  \end{qparts}
\end{question}

% ***************************** question 3 **************************

\begin{question}{20}
  \begin{qparts}

    \qpart[4]{A disease affects 1\% of people.  A test detects it in
      95\% of people who have it, but is also positive for 5\% of people
      who do not.  Someone tests positive.  What is the probability that
      they have the disease?
      \begin{answer}
        By Bayes' rule
        \begin{align*}
          \Prob{D\mid +} = \frac{0.95\times0.01}{0.95\times0.01 +
          0.05\times0.99} = \frac{0.0095}{0.0590} \approx 0.16.
        \end{align*}
        Most positives are false positives, because the disease is rare:
        the prior matters.
      \end{answer}}{\lines{5}}

    \qpart[4]{A coin with unknown probability $p$ of a head is tossed 10
      times, giving 7 heads.  Starting from the uniform prior
      $\mathrm{Beta}(p\mid1,1)$, what is the posterior?  Find its mean
      and its most probable value, and compare with the frequency
      $7/10$.
      \begin{answer}
        The posterior is $\mathrm{Beta}(p\mid1+7,1+3) =
        \mathrm{Beta}(p\mid8,4)$.  Its mean is $8/12\approx0.67$ and its
        mode $(8-1)/(8+4-2) = 0.7$, the frequency.  The mean is pulled
        slightly towards $\frac12$ by the prior.
      \end{answer}}{\lines{5}}

    \qpart[3]{Show that updating after each toss, using each posterior
      as the prior for the next toss, gives the same posterior as part
      (b), whatever the order of the heads and tails.
      \begin{answer}
        Each head adds 1 to the first parameter of the beta
        distribution and each tail adds 1 to the second.  After all ten
        tosses we have added 7 and 3 in some order, giving
        $\mathrm{Beta}(p\mid8,4)$ whatever the order.
      \end{answer}}{\lines{4}}

    \qpart[4]{The number of cars passing in an hour is Poisson with rate
      $\mu$, and the prior is $\mathrm{Gamma}(\mu\mid2,1)$.  On three
      days we count 3, 5 and 4 cars.  Using the update $a\to a+N$,
      $b\to b+1$ for each count, find the posterior and its mean and
      standard deviation ($a/b$ and $\sqrt{a}/b$).
      \begin{answer}
        $a = 2+3+5+4 = 14$ and $b = 1+3 = 4$, so the posterior is
        $\mathrm{Gamma}(\mu\mid14,4)$ with mean $3.5$ and standard
        deviation $\sqrt{14}/4\approx0.94$.
      \end{answer}}{\lines{4}}

    \qpart[5]{In Python, simulate 100 tosses of a coin with $p=0.7$ and
      update a $\mathrm{Beta}(1,1)$ prior after each toss.  Print the
      posterior mean and standard deviation,
      $\sqrt{ab/((a+b)^2(a+b+1))}$, after 10, 30 and 100 tosses.
      \begin{answer}
\begin{alltt}
rng = np.random.default\_rng(1)\\
a, b = 1, 1\\
for n, head in enumerate(rng.random(100) < 0.7, start=1):\\
\hspace*{2em}a += head; b += 1 - head\\
\hspace*{2em}if n in (10, 30, 100):\\
\hspace*{4em}print(n, a/(a+b), np.sqrt(a*b/((a+b)**2*(a+b+1))))
\end{alltt}
        With this seed the means are about $0.67$, $0.66$ and $0.67$ and
        the standard deviations $0.13$, $0.08$ and $0.05$: the
        uncertainty shrinks roughly as $1/\sqrt{n}$.
      \end{answer}}{\lines{6}}

  \end{qparts}
\end{question}

% ***************************** question 4 **************************

\begin{question}{25}
  This question is harder and optional.  When both the likelihood and
  the prior are normal, so is the posterior, and we only need to keep
  track of a mean and a variance (or a covariance matrix).  A helpful
  way to read the prior is as \emph{pseudo-data}: observations we
  pretend we have already seen.
  \begin{qparts}

    \qpart[5]{We measure an unknown quantity $\mu$ $n$ times, with
      independent errors: $X_i\sim\mathcal{N}(\mu,\sigma^2)$ with
      $\sigma$ known.  Our prior is $\mu\sim\mathcal{N}(\mu_0,s_0^2)$.
      Write down the log of the posterior, up to a constant, and show
      that it is quadratic in $\mu$.  Hence show the posterior is
      $\mathcal{N}(\mu_n, s_n^2)$ with
      \begin{align*}
        \frac{1}{s_n^2} = \frac{1}{s_0^2} + \frac{n}{\sigma^2},
        \qquad
        \mu_n = s_n^2\left(\frac{\mu_0}{s_0^2} +
        \frac{1}{\sigma^2}\sum_{i=1}^n X_i\right).
      \end{align*}
      (Hint: the log of $\mathcal{N}(\mu\mid m,s^2)$ is
      $-\mu^2/(2s^2) + \mu\,m/s^2 + $ a constant; match the
      coefficients of $\mu^2$ and $\mu$.)
      \begin{answer}
        \begin{align*}
          \log f(\mu\mid\mathcal{D}) &= -\frac{(\mu-\mu_0)^2}{2s_0^2}
          - \sum_{i=1}^n \frac{(X_i-\mu)^2}{2\sigma^2} + \text{const}\\
          &= -\frac{\mu^2}{2}\left(\frac{1}{s_0^2}+\frac{n}{\sigma^2}\right)
          + \mu\left(\frac{\mu_0}{s_0^2} + \frac{\sum_i X_i}{\sigma^2}\right)
          + \text{const}.
        \end{align*}
        Matching the coefficient of $\mu^2$ with $-1/(2s_n^2)$ and the
        coefficient of $\mu$ with $\mu_n/s_n^2$ gives the result.
        Inverse variances (\emph{precisions}) add.
      \end{answer}}{\lines{10}}

    \qpart[4]{Let $n_0 = \sigma^2/s_0^2$ and $\bar{X} = \frac1n\sum_i
      X_i$.  Show that
      $\mu_n = (n_0\,\mu_0 + n\,\bar{X})/(n_0+n)$ and
      $s_n^2 = \sigma^2/(n_0+n)$.  Interpret $n_0$.  What happens when
      $s_0\to\infty$, and when $n\to\infty$?
      \begin{answer}
        Multiplying top and bottom by $\sigma^2$,
        $1/s_n^2 = (n_0+n)/\sigma^2$ and
        $\mu_n = (n_0\mu_0 + n\bar{X})/(n_0+n)$.  The prior acts like
        $n_0$ extra measurements, all equal to $\mu_0$: a confident prior
        (small $s_0$) is worth many measurements.  As $s_0\to\infty$,
        $n_0\to0$, so $\mu_n = \bar{X}$ and $s_n = \sigma/\sqrt{n}$, the
        standard error from the lecture on experimental results.  As
        $n\to\infty$ the data swamp the prior.
      \end{answer}}{\lines{7}}

    \qpart[3]{Take $\mu_0=0$, $s_0=2$, $\sigma=1$ and measurements
      $1.2$, $0.8$, $1.5$, $0.9$.  Find $n_0$, $\mu_n$ and $s_n$.
      \begin{answer}
        $n_0 = 1/4$, $n=4$ and $\sum_i X_i = 4.4$, so
        $\mu_n = 4.4/4.25\approx1.04$ and $s_n = 1/\sqrt{4.25}\approx0.49$.
        The prior pulls the mean $\bar{X}=1.1$ slightly towards 0.
      \end{answer}}{\lines{4}}

    \qpart[6]{Now consider linear regression, $y_k = \bm{w}^\tr\bm{\phi}_k
      + \epsilon_k$ with $\epsilon_k\sim\mathcal{N}(0,\sigma^2)$, where
      the feature vectors $\bm{\phi}_k$ are the rows of a matrix
      $\mat{\Phi}$.  Take the prior $\bm{w}\sim\mathcal{N}(\bm{0},
      (\sigma^2/\eta)\,\mat{I})$: the weights are probably small.  Using
      the multivariate version of the hint,
      $\e{-\frac12\bm{w}^\tr\mat{A}\bm{w} + \bm{b}^\tr\bm{w}} \propto
      \mathcal{N}(\bm{w}\mid\mat{A}^{-1}\bm{b}, \mat{A}^{-1})$ (we show
      this in the lecture on Gaussian integrals), show that the posterior
      has precision $\mat{A} = (\mat{\Phi}^\tr\mat{\Phi}+\eta\,\mat{I})/\sigma^2$
      and mean $(\mat{\Phi}^\tr\mat{\Phi}+\eta\,\mat{I})^{-1}\mat{\Phi}^\tr\bm{y}$.
      What do you recognise?
      \begin{answer}
        \begin{align*}
          \log f(\bm{w}\mid\mathcal{D}) &= -\frac{\eta}{2\sigma^2}\len{\bm{w}}^2
          - \frac{1}{2\sigma^2}\len{\bm{y}-\mat{\Phi}\bm{w}}^2 + \text{const}\\
          &= -\frac12\bm{w}^\tr\frac{\mat{\Phi}^\tr\mat{\Phi}+\eta\mat{I}}{\sigma^2}\bm{w}
          + \bm{w}^\tr\frac{\mat{\Phi}^\tr\bm{y}}{\sigma^2} + \text{const},
        \end{align*}
        so $\mat{A} = (\mat{\Phi}^\tr\mat{\Phi}+\eta\mat{I})/\sigma^2$,
        $\bm{b} = \mat{\Phi}^\tr\bm{y}/\sigma^2$ and the mean is
        $\mat{A}^{-1}\bm{b} = (\mat{\Phi}^\tr\mat{\Phi}+\eta\mat{I})^{-1}\mat{\Phi}^\tr\bm{y}$:
        ridge regression.  The regulariser is the prior, and as in part
        (a) precisions add: each data point adds
        $\bm{\phi}_k\bm{\phi}_k^\tr/\sigma^2$.  The posterior also gives
        error bars on the weights, which ridge regression alone does not.
      \end{answer}}{\lines{10}}

    \qpart[7]{Fit a line $y = w_1 + w_2 x$ ($\bm{\phi}_k = (1,x_k)^\tr$)
      to the points $(0,1.1)$, $(1,2.9)$, $(2,5.2)$, $(3,6.8)$ with
      $\sigma=1$ and $\eta=0.5$.  In Python, update the posterior one
      point at a time, by adding to the precision and to
      $\bm{b}$, and print the posterior mean and the standard deviations
      of $w_1$ and $w_2$ after each point.  Check the final mean against
      ridge regression.
      \begin{answer}
\begin{alltt}
x = np.array([0.0, 1.0, 2.0, 3.0]); y = np.array([1.1, 2.9, 5.2, 6.8])\\
sig2, eta = 1.0, 0.5\\
A = (eta/sig2) * np.eye(2)\hspace*{2em}\# prior precision\\
b = np.zeros(2)\\
for xk, yk in zip(x, y):\\
\hspace*{2em}phi = np.array([1.0, xk])\\
\hspace*{2em}A = A + np.outer(phi, phi)/sig2\\
\hspace*{2em}b = b + phi*yk/sig2\\
\hspace*{2em}m = np.linalg.solve(A, b)\hspace*{2em}\# posterior mean\\
\hspace*{2em}sd = np.sqrt(np.diag(np.linalg.inv(A)))\\
\hspace*{2em}print(xk, m, sd)\\
Phi = np.column\_stack([np.ones(4), x])\\
print(np.linalg.solve(Phi.T @ Phi + eta*np.eye(2), Phi.T @ y))
\end{alltt}
        After the four points the posterior mean is $(1.02, 1.90)$,
        exactly the ridge solution, with standard deviations $0.70$ and
        $0.39$.  After one point (at $x=0$) the slope is still $0$ with
        its prior standard deviation $\sqrt2\approx1.41$: a single point
        says nothing about the slope.  Each new point shrinks the
        uncertainty.
      \end{answer}}{\lines{10}}

  \end{qparts}
\end{question}

\end{document}
